\documentclass[11pt, a4paper]{article}
\usepackage[utf8]{inputenc}
\usepackage{amsmath, amssymb}
\usepackage{graphicx}
\usepackage[margin=1in]{geometry}
\usepackage{tcolorbox}
\usepackage{fancyhdr}

\definecolor{UNIBBlue}{RGB}{0, 51, 102}

\pagestyle{fancy}
\fancyhf{}
\rhead{\textbf{Optimisasi untuk Teknik Elektro}}
\lhead{Problem Set - Minggu 5}
\cfoot{\thepage}

\begin{document}

\begin{center}
    \Large\color{UNIBBlue}\textbf{PROBLEM SET}\\
    \large\textbf{Minggu 5: Optimisasi Non-Linier (NLP) Tak Terkendala}
\end{center}

\vspace{0.5cm}

\textbf{Instruksi:} Kerjakan soal-soal di bawah ini dengan lengkap.

\section*{Soal 1: Metode Newton 1D (C3)}
Diberikan fungsi $f(x) = x^4 - 3x^3 + 2$.
Lakukan 2 iterasi Metode Newton secara manual dengan tebakan awal $x_0 = 3$. 
Tunjukkan nilai $f'(x)$, $f''(x)$, dan $x_{k+1}$ pada setiap iterasi.

\section*{Soal 2: Hessian Matriks (C4)}
Diberikan fungsi biaya dengan dua variabel $f(x_1, x_2) = 2x_1^2 + x_2^2 - x_1 x_2 - 4x_1$. \\
(a) Tentukan vektor gradien $\nabla f(x_1, x_2)$. \\
(b) Tentukan Matriks Hessian $H(x_1, x_2)$. \\
(c) Apakah Matriks Hessian tersebut \textit{Positive Definite}? Buktikan menggunakan nilai eigen (eigenvalues) atau determinan principal minor.

\section*{Soal 3: Komparasi dan Evaluasi (C5)}
Dalam aplikasi estimasi kanal (Channel Estimation) untuk telekomunikasi 5G, seorang insinyur ingin memperbarui bobot filter secara \textit{real-time} (setiap mikrodetik). Model error yang dioptimasi memiliki tingkat kelengkungan yang cukup konsisten, namun jumlah parameternya (variabel) mencapai $10^4$ (sepuluh ribu).

(C5) Metode mana yang akan Anda rekomendasikan untuk di-deploy ke perangkat keras (FPGA): Gradient Descent murni, Metode Newton, atau varian Quasi-Newton? Evaluasi kelebihan dan kekurangan ketiganya dalam konteks kebutuhan real-time dan ukuran variabel!

\section*{Soal 4: Sintesis Solusi Masalah Kegagalan (C6)}
Algoritma Gradient Descent yang diimplementasikan pada kontroler cerdas sebuah robot penjejak garis ternyata berosilasi (bergetar maju mundur) secara tidak terkendali di sekitar garis target dan tidak mau diam (konvergen).

(C6) Rancanglah dua buah modifikasi (strategi) secara matematis pada algoritma pembaruan (update rule) untuk memecahkan masalah osilasi ini. Jelaskan mekanisme kerja masing-masing strategi secara logis.

\end{document}
